\startmode[single]

\setvariables[meta]
  [title={Elektronenbeugunsröhre},
   subtitle={Versuchsprotokoll},
   date={2024-10-18},
   author={Niklas von Hirschfeld},
  ]

\setupinteraction[title={\getvariable{meta}{title}}]

\starttext

\startfrontmatter
\component c_titlepage
\component c_toc
\stopfrontmatter

\setuppagenumber[number=1]

\startbodymatter

\chapter{Versuchsprotokoll}

\stopmode

\starttabulate
\NC Name  \EQ \currentstructurename  \NC\NR
\NC Level \EQ \currentstructurelevel \NC\NR
\stoptabulate

In diesem Versuch werden Elektronen, mittels der Elektronenkanone, durch eine
vielzahl von übereinander gelegten und verdrehten {\em Graphitfolien} auf eine
Floreszenzschicht geschossen. 

Die Versuchsaufbau besteht aus:

\startitemize[packed]
\item Elektronenkanone
\item Graphitfolien
\item Glaskolben (mit Vakuum) und Floreszenzschicht
\stopitemize

Die Elektronen werden, mit der Elektronenkanone, durch die Graphitfolien
geschossen und stoßen dann auf die Floreszenzschicht, wodurch sie für uns
sichtbar und messbar werden.

% MP Grafik / Buffer fuer den Aufbau {{{
\startbuffer[fig:mp:wellen:beugungskanone:aufbau]
\startMPpage

u := 1cm;

z0=(-5 * u, 0.5 * u);
z1=(-5 * u, -0.5 * u);
z2=(-1 * u, 0.5 * u);
z3=(-1 * u, -0.5 * u);

z4=(-2u, 0.7u);
z5=(-2u, -0.7u);
z6=(2u, 0.7u);
z7=(2u, -0.7u);

z8=(3.5u, 0);


draw z0--z1--z3--z2--cycle;

label("Elektronenkanone", (.5[x0,x2],0.2u + .5[y0,y1]));
drawarrow (x0 + 0.2u, -0.2u + .5[y0,y1])--(x2 - 0.2u, -0.2u + .5[y0,y1]);

for a = -0.3 step 0.05 until 0.3:
	draw (a * u, 0.3 * u)--(a * u, -0.3 * u) withcolor .6white;
endfor;


draw (0, -1u)--(0, 0);

label.bot("Graphitfolien", (0, -1u));

draw (5.35u, -2u)--(5.35u, 0);

label.bot("Floreszenzschicht", (5.35u, -2u));

label("Vakuum", (2u, 0));

path c;
c = fullcircle scaled 4u shifted z8;

path d;
d = fullcircle scaled 6.5u shifted (x8 - 1.5u, y8);

z9=(z4--z6) intersectiontimes c;
z10=c intersectiontimes d;

% label.bot(decimal x9 & ", " & decimal y9, z6);
% label.bot(decimal x10 & ", " & decimal y10, z7);

draw subpath (y9, -y9) of c ;
draw subpath (y10, -y10) of d ;

draw z4--z6 cutafter c;
draw z5--z7 cutafter c;

\stopMPpage
\stopbuffer
% }}}

\placefigure[force][fig:mp:wellen:beugungskanone:aufbau]{Skizze des Aufbau}{
\typesetbuffer[fig:mp:wellen:beugungskanone:aufbau]
}


% MP Grafik / Buffer fuer die Beobachtung {{{
\startbuffer[fig:mp:wellen:beugungskanone:beobachtung]
\startMPpage
	pickup pencircle scaled 1pt;
    color green; green := (0, 1, 0);

    % Draw rings
    draw fullcircle scaled 30 withcolor green withpen pencircle scaled 2pt;  % inner most ring
    draw fullcircle scaled 50 withcolor green withpen pencircle scaled 2pt;  % second ring
    draw fullcircle scaled 60 withcolor green;  % third ring
    draw fullcircle scaled 95 withcolor green;  % outer most ring

    % Optionally, draw a small dot in the center to represent the electron spot
    fill fullcircle scaled 2 withcolor green;
\stopMPpage
\stopbuffer
% }}}

\margintext{
	\placefigure[force][fig:mp:wellen:beugungskanone:beobachtung]{Skizze der Beobachtung}{
	\typesetbuffer[fig:mp:wellen:beugungskanone:beobachtung]
	}
}

\startsectionlevel[title={Beobachtung}]

\starttabulate
\NC Name  \EQ \currentstructurename  \NC\NR
\NC Level \EQ \currentstructurelevel \NC\NR
\stoptabulate

% {{{
Nach anschalten der Elektronenkanone erscheinen vorne auf der Floreszenzschicht
ein grüner Punkt in der Mitte und Ringe um diesen herum, ca wie in
\in{Abbilung}[fig:mp:wellen:beugungskanone:beobachtung]. Bei änderung der
Beschleunigungsspannung  werden die Ringe enger und auch heller.

Gemessen wurde der Durchmesser der Ringe bei verschiedenen Beschleunigungsspannungen.

Auch lässt sich die Position der Ringe und des Punktes in der Mitte mittels
{\bf Magneten} beeinflussen. Wird ein Magnet oben an den Galskolben gehalten,
so verändert sich die Positionierung nach oben.
% }}}
\stopsectionlevel
\startsectionlevel[title={Auswertung}]
% {{{
Für die Auswertung wollen wir zum einen die Geschwindigkeit der Elektronen
wissen und auch die Wellenlänge berechnen.

\subsection{Berechnung der Geschwindigkeit}

Elektronen werden im homogenen elektrischen Feld (formal Kondensator)
beschleunigt.

\margintext{
\starttabulate[|m|p|]
\NC U_B \EQ Beschleunigungsspannung \NC\NR
\NC e \EQ Elementarladung \NC \NR
\NC m \EQ Masse des Elektron \NC \NR
\NC v \EQ Geschwindigkeit \NC \NR
\stoptabulate
} 

Im Feld besitzen sie $E_{el}=U_B \cdot e$ an Elektrischer Energie. Diese wird
(gemäß Annahme vollständig) in {\em kinetische} Energie ($E_{kin}=\frac{1}{2}m \cdot v^2$) umgewandelt. Diese beiden Energien können wir gleichstellen und nach $v$ (Geschwindigkeit) umstellen.

\startplaceformula[formel:wellen:verusch:beugungskanone:geschwindigkeit]
\startformula 
\startalign[n=3]
\NC E_{el}                  \NC = E_{kin}                   \NC \NR
\NC e \cdot U_B             \NC = \frac{1}{2} m \cdot v^2   \NC \quad \| \cdot 2 \| \div m \NR
\NC \frac{2eU_B}{m}         \NC = v^2                       \NC \quad \| \sqrt{} \NR
\NC \sqrt{\frac{2eU_B}{m}}  \NC = v                         \NC \NR
\stopalign
 \stopformula
\stopplaceformula

\startsectionlevel[title={Berechnung der Wellenlänge}]

Die Formel für die Wellenlänge ($\lambda$) ergibt sich wie folgt:

\startplaceformula[formel:beugungskanone:ringe]
\startformula
\tan \alpha = \frac{R}{l}
\stopformula
\stopplaceformula

Wobei $R$ der Radius und $l$ der Abstand der Graphitfolien zu dem Schrim
(\in{s. Abb.}[fig:mp:beugungskanone:ringe]). Für die Wellenlänge ergibt sich folgende Formel:

\startplaceformula[formel:beugungskanone:welenlaenge]
\startformula
\lambda \cdot n = a \sin(\alpha)
\stopformula
\stopplaceformula

Nun können wir die \in{Formel}[formel:beugungskanone:ringe] in die \in{Formel}[formel:beugungskanone:welenlaenge] einsetzen und erhalten:

\startplaceformula[formel:beugungskanone:ringe_eingesetzt]
\startformula
\lambda = \frac{a \cdot \sin(\text{arctan}\frac{R}{l})}{n}
\stopformula
\stopplaceformula

\starttabulate[|m|m|p|]
\NC a \NC \NC Abstand der Spalte (Gitterkonstante\sidenote{Durch die Graphitfolien, welche an sich verschiedene Gitterkonstanten aufweisen, wird hier mit zwei Verschiedene Gitterkonstanten ($d_1$ und $d_2$) gerechnet}) \NC\NR
\NC d_1 \EQ 123pm \NC Erste Gitterkonstante \NC \NR
\NC d_2 \EQ 213pm \NC Zweite Gitterkonstante \NC \NR
\NC R \EQ \frac{D}{2} \NC Radius der Ordnung \NC \NR
\NC L \EQ 130mm \NC Abstand zum Schirm \NC \NR
\stoptabulate

\startbuffer[fig:mp:beugungskanone:ringe]
\startMPpage

u := 2cm;

z0=-z3=(-2u, 0);
z1=(2u, 1u);
z2=(2u, -1u);

for a = -0.1 step 0.02 until 0.1:
	draw (x0 - a * u, 0.1 * u)--(x0 - a * u, -0.1 * u) withcolor .6white;
endfor;

path c;
c = fullcircle scaled 2.5u shifted z0;

draw c cutafter (z0 -- z1);

draw z0--z1;
draw z0--z2;

draw z0--z3;
draw z1--z3;

draw fullcircle yscaled 2u xscaled 1u shifted z3;

dotlabel("", z3);

% label.lft("R", 0.5[z1,z2]);
label.top(btex $\alpha$ etex, (-1*u, 0));
label.bot(btex $l$ etex, .5[z0,z3]);
label.rt(btex $R$ etex, .5[z1,z3]);

label.bot("Graphitfolie", (xpart z0, -0.1 * u));

\stopMPpage

\stopbuffer

\vtop{
\margintext{
	\starttabulate[|m|p|]
	\NC D \EQ Durchmesser Ring \NC\NR
	\NC R \EQ $\frac{D}{2}$ \NC\NR
	\NC l \EQ Abstand Graphit zum Leuchtschirm \NC\NR
	\stoptabulate
}
}

\placefigure[force][fig:mp:beugungskanone:ringe]{Ringe im Interferenzmuster}{
\typesetbuffer[fig:mp:beugungskanone:ringe]
}



\placetable[][table:wellen:versuch:beugungskanone:auswertung]{Auswertungen\sidenote{Diese Messungen sind nicht von uns selber. Es sind nur beispielhafte Zahlen.}}
{

\setupTABLE[r][1][style=bold,background=color,backgroundcolor=gray]
\bTABLE
% Header
\startCSV
U / V, $D_1$ / m, $D_2$ / m, v in m/s,	 $D_1$ zu $\lambda$, $D_2$ zu $\lambda$
\stopCSV
% Content
\processseparatedfile[CSV][assets/physik/wellenoptik/beugungskanone.csv]
\eTABLE
}

In der \in{Tabelle}[table:wellen:versuch:beugungskanone:auswertung] wurde, um
die Geschwindigkeit zu berechnen, die
\in{Formel}[formel:wellen:verusch:beugungskanone:geschwindigkeit] genutzt. Für
die Wellenlänge wurde die
\in{Formel}[formel:wellen:beugungskanone:wellenlaenge] genutzt. Da es zwei
verschiedenen Gitterkonstanten gibt, hat auch jede Ordnung {\bf zwei} Ringe mit jeweils unterschiedlichen Durchmessern ($D_1$ und $D_2$).

% MP Buffer fuer auswertung d1 und d2 {{{
\startbuffer[mp:wellen:versuch:beugungskanone:auswertung:d1]
\startMPpage
    % Define the size of the plot
    numeric u;
    u := 1cm;

    % Axes ranges
    numeric xmin, xmax, ymin, ymax;

    xmin := 1.50e-11;
    xmax := 3e-11;
    ymin := 2e7;
    ymax := 4.50e7;

    % Scaling factors to fit the plot into the canvas
    numeric xscale, yscale;
    xscale := 5u / (xmax - xmin);
    yscale := 3u / (ymax - ymin);

    % Function to transform (x, y) coordinates to the plot coordinates
    vardef transform(expr x, y) =
        ((x - xmin) * xscale, (y - ymin) * yscale)
    enddef;

    % Draw the axes
    drawarrow transform(xmin, ymin)--transform(xmax, ymin) withpen pencircle scaled 0.8pt;
    drawarrow transform(xmin, ymin)--transform(xmin, ymax) withpen pencircle scaled 0.8pt;

    % Draw the grid (optional)
    for i = 1.5e-11 step 5e-12 until 3e-11:
        draw transform(i, ymin)--transform(i, ymax) withpen pencircle scaled 0.3pt dashed evenly;
    endfor;
    for j = 2e7 step 5e6 until 4.5e7:
        draw transform(xmin, j)--transform(xmax, j) withpen pencircle scaled 0.3pt dashed evenly;
    endfor;

	label.bot("2e-11", transform(2e-11, ymin));
	label.bot("2.5e-11", transform(2.5e-11, ymin));
	label.bot("3e-11", transform(3e-11, ymin));

	label.lft("2.5e7", transform(xmin, 2.5e7));
	label.lft("3.5e7", transform(xmin, 3.5e7));
	label.lft("4.5e7", transform(xmin, 4.5e7));

    % Draw the points
	drawdot transform(2.68e-11, 2.97e7) withpen pencircle scaled 3pt; 
	drawdot transform(2.32e-11, 3.25e7) withpen pencircle scaled 3pt;
	drawdot transform(2.19e-11, 3.51e7) withpen pencircle scaled 3pt;
	drawdot transform(2.01e-11, 3.75e7) withpen pencircle scaled 3pt;
	drawdot transform(1.92e-11, 3.98e7) withpen pencircle scaled 3pt;
	drawdot transform(1.82e-11, 4.19e7) withpen pencircle scaled 3pt;

    % Label the axes
    label.top(btex $v \; \text{in} \frac{m}{s}$ etex, (0, 3u));
    label.urt(btex $\lambda \; \text{von} D_1$ etex, (5u, 0));

	% Define the function f(x)
	vardef f(expr x) =
	     6.65e-3 * (x**(-9.12e-1))
	enddef;

	% label(decimal f(2.68e-11), (2u, 2u));

	% Plot the function f(x) smoothly
    % pickup pencircle scaled 0.8pt;
    % drawoptions(withcolor red); % Set the color for the function curve
    path p;
    p = transform(1.75e-11, f(1.75e-11));
    % p = transform(1.5e-11, f(1.5e-11)) -- transform(2.8e-11, f(2.8e-11));
    % p = transform(2.68e-11, 2.97e7) -- transform(2.32e-11, 3.25e7);
    for x = 1.8e-11 step 5e-13 until 2.8e-11:
        p := p -- transform(x, f(x));
    endfor;
    draw p;

	label.rt(btex $f(x)=6.65 \text{e-03} \cdot x^{-9.12\mathrm{e}{-01}}$ etex, (0, -1u));
	label.rt(btex $R\approx 9,87 \text{e-01}$ etex, (0, -1.5u));
\stopMPpage
\stopbuffer


\startbuffer[mp:wellen:versuch:beugungskanone:auswertung:d2]
\startMPpage
    % Define the size of the plot
    numeric u;
    u := 1cm;

    % Axes ranges
    numeric xmin, xmax, ymin, ymax;

    xmin := 1.50e-11;
    xmax := 3e-11;
    ymin := 2e7;
    ymax := 4.50e7;

    % Scaling factors to fit the plot into the canvas
    numeric xscale, yscale;
    xscale := 5u / (xmax - xmin);
    yscale := 3u / (ymax - ymin);

    % Function to transform (x, y) coordinates to the plot coordinates
    vardef transform(expr x, y) =
        ((x - xmin) * xscale, (y - ymin) * yscale)
    enddef;

    % Draw the axes
    drawarrow transform(xmin, ymin)--transform(xmax, ymin) withpen pencircle scaled 0.8pt;
    drawarrow transform(xmin, ymin)--transform(xmin, ymax) withpen pencircle scaled 0.8pt;

    % Draw the grid (optional)
    for i = 1.5e-11 step 5e-12 until 3e-11:
        draw transform(i, ymin)--transform(i, ymax) withpen pencircle scaled 0.3pt dashed evenly;
    endfor;
    for j = 2e7 step 5e6 until 4.5e7:
        draw transform(xmin, j)--transform(xmax, j) withpen pencircle scaled 0.3pt dashed evenly;
    endfor;


	label.bot("2e-11", transform(2e-11, ymin));
	label.bot("2.5e-11", transform(2.5e-11, ymin));
	label.bot("3e-11", transform(3e-11, ymin));

	label.lft("2.5e7", transform(xmin, 2.5e7));
	label.lft("3.5e7", transform(xmin, 3.5e7));
	label.lft("4.5e7", transform(xmin, 4.5e7));

    % Draw the points
	drawdot transform(2.60e-11, 2.97e7) withpen pencircle scaled 3pt; 
	drawdot transform(2.36e-11, 3.25e7) withpen pencircle scaled 3pt;
	drawdot transform(2.12e-11, 3.51e7) withpen pencircle scaled 3pt;
	drawdot transform(2.04e-11, 3.75e7) withpen pencircle scaled 3pt;
	drawdot transform(1.88e-11, 3.98e7) withpen pencircle scaled 3pt;
	drawdot transform(1.80e-11, 4.19e7) withpen pencircle scaled 3pt;

    % Label the axes
    label.top(btex $v \; \text{in} \frac{m}{s}$ etex, (0, 3u));
    label.urt(btex $\lambda \; \text{von} D_1$ etex, (5u, 0));

	% Define the function f(x)
	vardef f(expr x) =
	     4.97e-3 * (x**(-9.23e-1))
	enddef;

	% label(decimal f(2.68e-11), (2u, 2u));

	% Plot the function f(x) smoothly
    % pickup pencircle scaled 0.8pt;
    % drawoptions(withcolor red); % Set the color for the function curve
    path p;
    p = transform(1.75e-11, f(1.75e-11));
    % p = transform(1.5e-11, f(1.5e-11)) -- transform(2.8e-11, f(2.8e-11));
    % p = transform(2.68e-11, 2.97e7) -- transform(2.32e-11, 3.25e7);
    for x = 1.8e-11 step 5e-13 until 2.8e-11:
        p := p -- transform(x, f(x));
    endfor;
    draw p;

	label.rt(btex $f(x)=6.65 \mathrm{e}{-03} \cdot x^{-9.12\mathrm{e}{-01}}$ etex, (0, -1u));
	label.rt(btex $R\approx 9,93 \mathrm{e}{-01}$ etex, (0, -1.5u));
\stopMPpage
\stopbuffer
% }}}

\startframedtext[width=\dimexpr\textwidth+\rightmarginwidth+\rightmargindistance\relax,frame=off] % This spans textwidth + margin width

\placefigure[force][mp:wellen:versuch:beugungskanone:auswertung:d1]{$v$-$\lambda$-Diagram und Regression für die Durchmesser $D_1$ und $D_2$}{
    \startcombination[2*1]
        {\typesetbuffer[mp:wellen:versuch:beugungskanone:auswertung:d1]}{Für $D_1$}
        {\typesetbuffer[mp:wellen:versuch:beugungskanone:auswertung:d2]}{Für $D_2$}
    \stopcombination
}

\stopframedtext
% }}}

Mögliches Runden:

\startformula
\lambda = \frac{0,0067^{\frac{m^2}{s}}}{v} 
\stopformula

% \startformula
% \startalign
%
% \NC \[\frac{1}{v}\] \NC =\frac{1}{\frac{m}{s}}=\frac{s}{m} \NC \qquad \[\lambda=m \NC \NR
% \NC \[\text{Vorfaktor}\] \NC = \frac{m^2}{s} \NC \NC \NR
%
% \stopalign
% \stopformula

\stopsectionlevel
\stopsectionlevel
\startsectionlevel[title={Fazit}]
\stopsectionlevel

\startmode[single]

\stopbodymatter

\startbackmatter
% \component c_definitions
% \component c_bibliography
\stopbackmatter
\stoptext

\stopmode
